\section{Тест производительности}
Для проверки производительности решение сравнивается с наивным алгоритмом, который для каждого запроса перебирает все отрезки
и проверяет условие $l \le x \le r$ и $h > y$. Его сложность равна $O(nm)$, а у реализованного решения $O((n + m)\log n)$.
Цель теста --- убедиться, что теоретические оценки подтверждаются на практике: при увеличении размера входных данных
время работы перебора должно расти квадратично, а время работы персистентного дерева отрезков --- почти линейно.

\subsection*{Условия эксперимента}
Размеры входных данных брались равными $n = m$ и удваивались от $1000$ до $64000$. Левые и правые границы отрезков,
высоты и координаты точек выбирались случайно из диапазона $[-10^9, 10^9]$ при помощи генератора со значением
\texttt{seed} $= 1$, поэтому почти все высоты различны, и дерево не вырождается. Обе программы компилировались
командой \texttt{g++ -O2 -std=c++17}. Для каждого размера решение запускалось $3$ раза, а наивный алгоритм --- $3$ раза
при $n < 50000$ и один раз при $n \ge 50000$; в таблице приведена медиана времени работы процесса.
Время включает чтение входных данных, вывод результата направлялся в \texttt{/dev/null}.

\subsection*{Проверка корректности}
Перед замерами ответы решения сравнивались с ответами наивного алгоритма на $200$ случайных малых
тестах ($n, m \le 8$, координаты по модулю не более $6$, поэтому часто встречаются совпадающие координаты и высоты). Расхождений нет.

\begin{alltt}
$ g++ -O2 -std=c++17 -o sol main1.cpp
$ g++ -O2 -std=c++17 -o brute brute.cpp
$ ./check.sh
OK: 200 малых тестов совпали
\end{alltt}

\subsection*{Результаты}
\begin{alltt}
$ python3 bench.py 1000 2000 4000 8000 16000 32000 64000
     n=m     sol, c   brute, c
    1000     0.0082     0.0093
    2000     0.0124     0.0202
    4000     0.0214     0.0595
    8000     0.0405     0.2062
   16000     0.0782     0.7886
   32000     0.1953     2.9173
   64000     0.3651    11.0841
\end{alltt}

В таблице ниже для каждого удвоения $n$ указано, во сколько раз выросло время по сравнению с предыдущей строкой.
При линейной зависимости ожидается значение около $2$, при квадратичной --- около $4$.

\begin{center}
\begin{tabular}{|r|c|c|c|c|}
\hline
\rowcolor{lightgray}
$n = m$ & Решение, с & Рост & Перебор, с & Рост \\
\hline
1000  & 0,0082 & ---  & 0,0093  & ---  \\
\hline
2000  & 0,0124 & 1,51 & 0,0202  & 2,17 \\
\hline
4000  & 0,0214 & 1,73 & 0,0595  & 2,95 \\
\hline
8000  & 0,0405 & 1,89 & 0,2062  & 3,47 \\
\hline
16000 & 0,0782 & 1,93 & 0,7886  & 3,82 \\
\hline
32000 & 0,1953 & 2,50 & 2,9173  & 3,70 \\
\hline
64000 & 0,3651 & 1,87 & 11,0841 & 3,80 \\
\hline
\end{tabular}
\end{center}

\subsection*{Выводы по результатам}
\begin{itemize}
\item \textbf{Наивный перебор} при каждом удвоении $n$ замедляется в $3{,}5$--$3{,}8$ раза (при $n \ge 8000$),
то есть примерно в $4$ раза, как и должно быть при сложности $O(n^2)$ (при $n = m$).
\item \textbf{Персистентное дерево отрезков} при каждом удвоении $n$ замедляется примерно в $1{,}9$ раза
(кроме выброса $2{,}5$ при $n = 32000$, который, вероятно, объясняется шумом измерения), то есть примерно в $2$ раза с
небольшой добавкой от множителя $\log n$. Это согласуется с оценкой $O((n + m)\log n)$.
\item При малых $n$ (порядка тысяч) времена обоих алгоритмов сопоставимы, а рост времени отличается от теоретического: в этих
замерах заметную долю составляют запуск процесса и ввод-вывод, а не сами вычисления. По мере роста $n$ эти накладные расходы
перестают влиять, и картина становится близкой к теоретической.
\item При $n = m = 64000$ решение работает за $0{,}37$ с, а наивный перебор --- за $11{,}1$ с, то есть примерно в $30$ раз
медленнее. Так как рост времени перебора квадратичный, а решения --- почти линейный, разрыв увеличивается с
ростом $n$ (при $n = 8000$ разница около $5$ раз, при $n = 16000$ --- около $10$ раз).
\end{itemize}

Таким образом, эксперимент подтвердил оценки сложности: $O(nm)$ для наивного алгоритма и $O((n + m)\log n)$ для персистентного
дерева отрезков.

\pagebreak